% Example 3 (fragmentation, interleaved): easy and hard steps alternate, so AI-friendly
% steps are never adjacent and no multi-step chain can form. Each easy step is augmented
% on its own (a length-1 chain) and each hard step is performed manually. Colors follow
% Figures 1-2 (manual=gray, automated=green, augmented=orange).
\begin{center}
\begin{adjustbox}{max width=\linewidth}
\begin{tikzpicture}[>=latex,
  cbox/.style={rectangle, rounded corners, draw, minimum width=2.0cm, minimum height=0.95cm, align=center, font=\small},
  manual/.style={fill=gray!10},
  automated/.style={fill=green!30},
  augmented/.style={fill=orange!30},
  sub/.style={draw=none, font=\scriptsize, text=black!70}]

  \node[cbox, augmented] (s1) {Step 1\\(Augmented)};
  \node[cbox, manual, right=0.7cm of s1] (s2) {Step 2\\(Manual)};
  \node[cbox, augmented, right=0.7cm of s2] (s3) {Step 3\\(Augmented)};
  \node[draw=none, right=0.45cm of s3] (dots) {$\cdots$};
  \node[cbox, augmented, right=0.45cm of dots] (sm1) {Step $m{-}1$\\(Augmented)};
  \node[cbox, manual, right=0.7cm of sm1] (sm) {Step $m$\\(Manual)};

  \node[sub, below=1pt of s1] {easy};
  \node[sub, below=1pt of s2] {hard};
  \node[sub, below=1pt of s3] {easy};
  \node[sub, below=1pt of sm1] {easy};
  \node[sub, below=1pt of sm] {hard};

  \draw[->,thick] (s1)--(s2);
  \draw[->,thick] (s2)--(s3);
  \draw[->,thick] (s3)--(dots);
  \draw[->,thick] (dots)--(sm1);
  \draw[->,thick] (sm1)--(sm);
\end{tikzpicture}
\end{adjustbox}
\end{center}
